% SPDX-License-Identifier: Apache-2.0
%
% The second TxHDL article: embedding the language in Rust.
% Built by //docs:embedding. See docs/document-build-plan.md.
\documentclass[journal,9pt]{IEEEtran}

% SPDX-License-Identifier: Apache-2.0
%
% The house preamble, shared by every document in this directory. Loaded
% after \documentclass, before \title. Keeping it in one file is what
% stops two articles drifting apart in font, listing style or figure
% style.
% Computer Modern. IEEEtran selects Times (ptm) by default, so the face has
% to be chosen explicitly.
%
% Latin Modern is the Computer Modern variety used here, and the choice is
% forced rather than aesthetic. Under T1 encoding, plain `cmr` has no Type 1
% outlines unless cm-super is installed, and it is not in the pinned
% distribution, so pdflatex falls back to EC bitmaps and every glyph in the
% PDF becomes a Type 3 font. That was measured: the first build produced a
% document with 10 Type 3 fonts and no Type 1. Latin Modern is Computer
% Modern redrawn with a full T1 repertoire and real .pfb outlines, so the
% same design comes out scalable.
\usepackage[T1]{fontenc}
\usepackage[utf8]{inputenc}
\usepackage{lmodern}
% microtype is not in the pinned TeX distribution. emergencystretch alone
% keeps the two-column measure from producing overfull lines.
\setlength{\emergencystretch}{3em}

\usepackage{amsmath}
\usepackage{amssymb}
\usepackage{amsthm}
\usepackage{booktabs}
\usepackage{array}
% enumitem is not in the pinned TeX distribution; the list spacing below
% does what \setlist would have done.
\usepackage{float}
\usepackage{xcolor}
\usepackage{listings}
\usepackage{tikz}
\usetikzlibrary{arrows.meta,positioning,fit,backgrounds,calc}
\usepackage[hidelinks,bookmarks=true]{hyperref}

% Numbered, definition-style box for a rule the reader is meant to apply.
\theoremstyle{definition}
\newtheorem{rulebox}{Rule}
\newtheorem{example}{Example}

% Unnumbered asides.
\theoremstyle{remark}
\newtheorem*{tip}{Tip}
\newtheorem*{pitfall}{Pitfall}

% Tighter lists, in place of enumitem's \setlist.
\setlength{\topsep}{2pt}
\setlength{\itemsep}{1pt}
\setlength{\parsep}{0pt}

\newcommand{\code}[1]{\texttt{#1}}
\newcommand{\term}[1]{\emph{#1}}

% Syntax highlighting. Four roles, distinguishable in colour and still
% distinguishable in greyscale, because an IEEE article gets printed: the
% keyword is the only bold role, the comment is the only italic one, and the
% two remaining roles differ in lightness as well as in hue.
\definecolor{kwblue}{RGB}{20,60,150}
\definecolor{cmtgreen}{RGB}{90,120,90}
\definecolor{strred}{RGB}{150,50,40}
\definecolor{numgray}{gray}{0.45}
\definecolor{framegray}{gray}{0.70}
\definecolor{bgshade}{gray}{0.97}
% A darker fill for figure cells. The listing background has to stay very
% light to keep code readable; a figure cell has to be clearly shaded or the
% distinction it draws is invisible in print.
\definecolor{cellshade}{gray}{0.90}

% The listing size is the document's choice, because it depends on the
% column: a two-column article sets `\listingsize` to 6pt and keeps its
% listings within 70 columns, a one-column document keeps the body size
% and includes source files at 80. Lines are never broken; a line that does not fit
% overflows the frame, which is visible, rather than wrapping, which is
% not.
\providecommand{\listingsize}{\scriptsize}

% `'` is deliberately not a string delimiter: the language uses it for loop
% labels, as Rust does, so treating it as a quote would swallow the rest of
% the line.
\lstdefinestyle{txhdl}{
  basicstyle=\ttfamily\listingsize,
  keywordstyle=\bfseries\color{kwblue},
  commentstyle=\itshape\color{cmtgreen},
  stringstyle=\color{strred},
  numberstyle=\tiny\color{numgray},
  morekeywords={mod,use,pub,const,type,struct,enum,trait,impl,fn,async,let,mut,
    match,if,else,loop,while,for,in,break,continue,return,as,await,
    module,bus,channel,transaction,protocol,pipeline,flow,seq,config,
    port,stage,pipe,instance,connect,bind,tag,domain,blackbox,foreign,
    property,role,signal,handler,true,false,
    sequence,interface,modport,implements,package,import,var,wire,func,proc,
    case,default,next,fork,branch,join,state,fsm},
  morecomment=[l]{//},
  morestring=[b]",
  showstringspaces=false,
  columns=fullflexible,
  keepspaces=true,
  breaklines=false,
  % Framed, so a listing is bounded rather than floating in the column.
  frame=single,
  rulecolor=\color{framegray},
  framesep=3pt,
  backgroundcolor=\color{bgshade},
  xleftmargin=3pt,
  xrightmargin=3pt,
  aboveskip=6pt,
  belowskip=6pt,
  % Every listing is numbered and captioned, so it can be referred to by
  % number. `caption` without `float` numbers the listing and leaves it
  % where it was written, which is what a listing in running prose needs.
  captionpos=b,
  abovecaptionskip=2pt,
  belowcaptionskip=6pt,
}

% Figures drawn as text. Framed the same way, and with the highlighting
% switched off, because a box-and-arrow diagram is not code and half its
% words turning blue reads as an accident. Setting a style adds keys rather
% than clearing the ones already set, so the three roles are neutralised by
% hand rather than by leaving them out.
\lstdefinestyle{ascii}{
  basicstyle=\ttfamily\listingsize,
  keywordstyle=\ttfamily,
  commentstyle=\ttfamily,
  stringstyle=\ttfamily,
  columns=fullflexible,
  keepspaces=true,
  frame=single,
  rulecolor=\color{framegray},
  framesep=3pt,
  backgroundcolor=\color{bgshade},
  xleftmargin=3pt,
  xrightmargin=3pt,
  aboveskip=6pt,
  belowskip=6pt,
}
\lstset{style=txhdl}

% House TikZ styles. `shorten` lives on the arrow styles rather than on each
% arrow, so every arrow in the document gets the gap, including ones added
% later. TikZ clips a path at a node's border, so without it an arrowhead
% butts against the box with no gap at all. Keep `node distance` well above
% twice the shorten, or the arrow shrinks to nothing.
\tikzset{
  box/.style={draw,rounded corners=2pt,align=center,inner sep=4pt,
              font=\footnotesize},
  boxf/.style={box,fill=bgshade},
  lbl/.style={font=\scriptsize\itshape,align=center},
  fl/.style={-{Stealth[length=4pt]},shorten >=2.5pt,shorten <=2.5pt},
  fld/.style={fl,dashed},
}



\InputIfFileExists{buildstamp.tex}{}{}
\providecommand{\buildstamp}{Draft build (outside the Bazel flow).}

\title{Deleting a Language: Embedding TxHDL in Rust}

\author{Filip Filmar and Dragisa Jankovic%
\thanks{This article follows the TxHDL merge reported in the companion
paper. It asks what remains of the language once every construct Rust
already provides is removed. Every claim about Rust in it was compiled;
the probes are in the project repository under
\code{experiments/rust\_embedding}.}%
\thanks{\buildstamp}}

\markboth{TxHDL Embedded in Rust}{Filmar and Jankovic: Deleting a Language}

\begin{document}
\maketitle

\begin{abstract}
A companion paper merged two hardware description languages into one and
gave the result a Rust-shaped syntax. This paper asks the next question: if
the syntax is already Rust, why is there a separate language at all. We
remove every construct that Rust provides and keep only what is left. Of
the twenty non-Rust constructs in the merged specification, seventeen go.
Two results are worth the exercise on their own. The rule that a timeless
body may not count cycles needs no enforcement, because a cycle count is
spelled \code{.await} and a plain \code{fn} has none. And a pipeline, a
sequence and a set of parallel processes are one construct, an \code{async
fn}, differing only in how they are driven; the live-range analysis that a
\code{pipe} declaration existed to feed is what an async state machine does
already. Every claim here was compiled rather than argued, across sixteen probes,
and four of those compilations changed the design. Arbitrary-width
arithmetic needs nightly Rust. A unit's parallel processes cannot take
\code{\&mut self}. A \code{macro\_rules!} conditional cannot reach inside
a block, which is the strongest argument for elaborating through a
procedural macro. And one driver per wire turns out to be free: a signal
yields its two ends separately, the driving end is not \code{Clone}, so
the oldest error message in hardware design arrives as a move error. One
wall remains, and it is the reason the merged language was designed with
its own parser: \code{if} on a signal cannot work, because a signal has no
value while the program runs. Half of the result is Chisel's architecture
in Rust, and is presented as such. The other half, sequencing through
\code{async}, is the contribution, its nearest ancestors are Calyx and
Filament rather than Chisel, and it is the half that has not yet been
lowered to hardware. The paper closes by stating the one experiment that
would settle whether it can be.
\end{abstract}

\begin{IEEEkeywords}
Hardware description languages, embedded domain-specific languages, Rust,
high-level synthesis, async, language design.
\end{IEEEkeywords}

{\footnotesize\tableofcontents}

% SPDX-License-Identifier: Apache-2.0
\section{Introduction}
\label{sec:e-intro}

\IEEEPARstart{A}{} companion paper~\cite{merge} took two hardware
description languages that had been developed independently for one
project, showed that they did not compete, and merged them. LHDL described
dataflow and left time to the compiler. TxHDL described transactions over
buses and let a module block for as many cycles as a protocol took. Four
questions needed a decision about meaning, and the analysis behind
them~\cite{unification} resolved all four.

The merged language was then given a surface syntax biased toward
Rust~\cite{syntax}, construct by construct, because Rust is familiar and
because following it removed six defects that the two sources had rather
than restating them.

This paper asks the question that arrangement invites. If the syntax is
already Rust, why is there a separate language at all?

\subsection{The experiment}

Take the merged specification and delete every construct that Rust already
provides. Rename every construct whose name collides with a Rust concept,
so that nothing means two things. Keep what is left, and see how small it
is.

The answer is that seventeen of the twenty non-Rust constructs go, and what
remains is a library rather than a keyword list.

Two of the deletions pay for the exercise on their own, and
Sections~\ref{sec:e-time} and~\ref{sec:e-pipeline} state them. A third
result is negative and matters as much: Section~\ref{sec:e-walls} gives the
one thing an embedded language cannot do, which is the thing a separate
parser was buying.

\subsection{Everything here was compiled}

The first draft of this work argued about what Rust would accept. That was
replaced. A Rust toolchain is now part of the repository's hermetic build,
and every claim in this paper is a probe that was compiled, one per
question. Section~\ref{sec:e-probes} reports what each one answered, and
keeps the two that failed, because a failure is an answer.

Two of those compilations changed the design rather than confirming it.
Arbitrary-width arithmetic does not work on stable Rust. A unit's parallel
processes cannot take a mutable borrow of the unit. Neither was visible
from the language reference; both were visible in ten lines of code.

% SPDX-License-Identifier: Apache-2.0
\section{The Rule, and What It Leaves}
\label{sec:e-rule}

\begin{rulebox}[The reduction]
Delete every construct Rust already provides. Rename every construct whose
name collides with a Rust concept. What remains is the library.
\end{rulebox}

The renaming clause does one job in practice. A hardware \term{module} and
a Rust \code{mod} are different things with one name, so the hardware one
becomes a \term{unit} and \code{mod} keeps its meaning as a namespace.

\subsection{The vocabulary that remains}

Two Rust modules hold everything the language still needs.

\begin{lstlisting}
crate::types::Transaction   // data that moves between units
crate::types::Tag           // a domain and its policy

crate::comp::Module         // a unit
crate::comp::Bus            // a bundle of channels
crate::comp::Chan           // one channel of a bus
crate::comp::Signal         // one wire
crate::comp::Port           // an attachment to a bus, in a role
crate::comp::Role           // the type-level tag for a modport
\end{lstlisting}

Everything else is a plain Rust struct, trait or function.
Table~\ref{tab:reduction} states the mapping.

\begin{table}[H]
\caption{What each construct becomes. Seventeen of twenty are deleted
outright.}
\label{tab:reduction}
\centering
\footnotesize
\begin{tabular}{@{}ll@{}}
\toprule
Was & Is now \\
\midrule
\code{transaction T}   & \code{struct T} plus \code{impl Transaction} \\
\code{module M}        & \code{struct M}, a unit, plus \code{impl Module} \\
\code{bus B}           & \code{struct B} plus \code{impl Bus} \\
\code{interface I}     & \code{struct I} of \code{Signal} fields \\
\code{modport R}       & \code{trait R}, implemented for the interface \\
\code{channel c: T}    & a field of type \code{Chan<T>} \\
\code{port p}          & a field of type \code{Port<B, R>} \\
\code{flow f}          & \code{fn f} \\
\code{seq s}           & \code{async fn} \\
\code{pipeline stage pipe} & \code{async fn} and \code{.await} \\
\code{tag T}           & \code{struct T} plus \code{impl Tag}, as a parameter \\
\code{config bind}     & a \code{type} alias \\
\code{domain D}        & a type parameter \\
\code{assert assume cover} & Rust's own macros \\
\bottomrule
\end{tabular}
\end{table}

\subsection{Transactions, interfaces and modports}

A transaction is a struct that implements a marker trait. Nothing else
about it is special.

\begin{lstlisting}
pub struct MacRequest { pub addr: u32, pub count: u8 }
impl Transaction for MacRequest {}
\end{lstlisting}

An interface splits into two Rust items, and that split fixes a defect the
earlier analysis reported~\cite{unification}. LHDL's \code{interface} did
two jobs at once, a bundle of wires and a socket an implementation fills,
which is why its grammar and its examples never agreed. Here the bundle is
a struct and the role is a trait.

\begin{lstlisting}
// The wires. No direction on any of them.
pub struct Wishbone {
    pub adr: Signal<u32>,
    pub ack: Signal<bool>,
}

// A role. Direction lives here, not in the struct.
pub trait Initiator {
    fn drive_adr(&mut self, v: u32);
    fn read_ack(&self) -> bool;
}

impl Initiator for Wishbone { /* ... */ }
\end{lstlisting}

A unit may then be generic over the role it attaches in rather than over a
concrete bus, which is what late binding needs and
Section~\ref{sec:e-binding} develops.

% SPDX-License-Identifier: Apache-2.0
\section{Values}
\label{sec:e-values}

Rust's integers are the wrong types for hardware, and the reason is not
width. It is that \code{u32} has no way to say that nobody has driven this
wire yet.

\code{crate::types} therefore defines its own, paralleling the primitives
rather than reusing them.

\subsection{Two values and nine}

\code{Bit} is two valued, and it is what synthesis maps. Its default is
\code{Zero}, because a synthesised register leaves reset at a defined
value.

\code{Logic} is nine valued, in the order IEEE 1164 gives
\code{std\_ulogic}. Its default is \code{U}, uninitialised, and that
difference is the whole reason the type exists.

\begin{lstlisting}[caption={The two value types. An undriven wire is \code{U}, not zero.},label={lst:bitlogic}]
pub enum Bit { Zero, One }

pub enum Logic {
    U,        // uninitialised
    X,        // forcing unknown
    Zero, One,
    Z,        // high impedance
    W,        // weak unknown
    L, H,     // weak 0, weak 1
    DontCare,
}
\end{lstlisting}

Narrowing is fallible, and that is the point. \code{Logic::to\_bit}
returns nothing for \code{U}, \code{X} or \code{Z}, so a design that turns
an unknown into a number has to say so rather than do it quietly.

\subsection{Resolution, not a rule of our own}

Two drivers on one wire are resolved by the IEEE 1164 table, which the
library implements rather than paraphrases.

\begin{lstlisting}[caption={Resolution. Contention gives \code{X}; a pull-up against high impedance gives \code{H}.},label={lst:resolve}]
Logic::One.resolve(Logic::Zero)   // X
Logic::Z.resolve(Logic::H)        // H
\end{lstlisting}

Reusing the standard table costs nothing and means an engineer who knows
VHDL already knows what this does.

\subsection{Numbers}

\code{U<N>} and \code{I<N>} are the numeric vectors, and
\code{logic::Vec<N>} is a vector of \code{Logic} for when unknowns must
propagate. It lives in its own module so that the type is just
\code{Vec}: \code{logic::Vec<8>} says what it is without repeating the
word, and \code{std::vec::Vec} is untouched, because nothing is imported
unqualified. Converting between them is where the two worlds meet:
\code{logic::Vec::to\_u} returns nothing unless every bit is defined.

\begin{lstlisting}[caption={An undriven vector is unknown, and says so.},label={lst:undriven}]
let v = logic::Vec::<8>::default();   // eight U's
assert!(v.to_u().is_none());
\end{lstlisting}

Same-width arithmetic, resizing to a stated width and slicing all compile
on stable, because each result width is a standalone const parameter. So
does a multiply whose result width is \emph{stated}, \code{a.mul::<64>(b)},
because \code{64} is standalone too. What does not compile is a multiply
whose result width is \emph{computed}, \code{U<\{A + B\}>}, for the
reason Section~\ref{sec:e-probes} gives: that is the one thing stable
Rust will not accept, so the width is written by the caller rather than
derived.

A number is written as a number. \code{U<N>} converts from every unsigned
primitive and from \code{i32}, and its operators and every drive take
anything that converts, so \code{acc + 1}, \code{n == 0} and
\code{reg.set(0)} say what they mean. The operators are Rust's,
\code{+}, \code{-}, \code{\&}, \code{|}, \code{\^{}}, \code{!},
\code{<<}, \code{>>} and the compares, through the \code{std::ops}
traits; a compare yields a \code{bool}, since \code{PartialEq} allows
nothing else, and a \code{bool} converts to a \code{Bit} and back, so
a condition is either. \code{i32} is there because an
integer literal with no other constraint is an \code{i32}, and a literal
is the common case; a negative one is a bug, and debug builds say so.

% SPDX-License-Identifier: Apache-2.0
\section{The Model of Time Enforces Itself}
\label{sec:e-time}

The merged specification spends a section on one rule. A timeless body may
not count cycles; a sequenced body counts them exactly as written. It gives
the two bodies two keywords, \code{flow} and \code{seq}, and states that a
cycle count inside a \code{flow} is a compile error.

Embedded in Rust, that rule needs no enforcement, because it cannot be
broken.

A cycle count is spelled \code{.await}. A plain \code{fn} is not
\code{async}, so it has no \code{.await} to write, so it cannot state a
cycle count. The compiler rejects the attempt already, with its own error
about awaiting outside an async context, and the specification says
nothing.

\begin{lstlisting}[caption={The model of time, enforced by the type system. The first function has no \code{.await} to write, so it cannot state a cycle count.},label={lst:time}]
// Timeless: combinational. There is no .await to write.
fn narrow(v: U<64>, low: bool) -> U<32> {
    mux(low, low_half(v), high_half(v))
}

// Sequenced. Counts cycles exactly as written.
async fn main(&mut self) {
    loop {
        let req = self.host.request.recv().await;
        cycles(BAUD_TICKS).await;
    }
}
\end{lstlisting}

Two constructs are deleted, and a rule that needed a paragraph of
specification and a compiler check becomes a fact about the type system.

The division is not between cheap operations and expensive ones. It is
between operations that finish inside a cycle and operations that do not.
A multiply is \code{async} for that reason, and
Section~\ref{sec:e-pipeline} shows what follows from it.

This is the clearest case of the reduction paying for itself, and it is
worth saying why it works rather than treating it as luck. The original
rule and Rust's rule are the same rule. LHDL's eighth thesis~\cite{theses}
argued that a designer should state intent and leave the mapping to time to
the implementation. Rust's \code{async} exists to mark exactly the place
where a computation gives up control and resumes later. Both are drawing
the same line, so one of them can be deleted.

\input{embedding_sections/05_pipeline}
% SPDX-License-Identifier: Apache-2.0
\section{Units, Processes and Parallelism}
\label{sec:e-units}

A unit is a struct that implements \code{crate::comp::Module}. The trait
has one method.

\begin{lstlisting}[caption={The unit trait. Inputs and outputs are type parameters, so one struct may implement it more than once.},label={lst:module}]
#[allow(async_fn_in_trait)]
pub trait Module<In, Out> {
    async fn run(&mut self, inputs: In, outputs: Out);
}
\end{lstlisting}

Inputs and outputs are type parameters rather than associated types, and
that choice buys something concrete: one struct may implement the trait
more than once, with different shapes, so a unit that works on 32-bit and
64-bit streams is one implementation and not two.

\subsection{Parallelism is stated, not implied}

Synthesis calls \code{run}. A unit with more than one concurrent process
starts one async function per process and joins them.

\begin{lstlisting}[caption={Parallelism written down. A unit with one process writes no join.},label={lst:join}]
impl Module<Ins, Outs> for DualPort {
    async fn run(&mut self, _i: Ins, _o: Outs) {
        join2(self.port_a(), self.port_b()).await;
    }
}
\end{lstlisting}

This replaces a rule with a construct. TxHDL said that several sequence
blocks in one module run concurrently, which was a property of having
declared them. It also contradicted itself: the specification stated that a
module holds one sequence, while three of its own modules declared
two~\cite{unification}. Here there is no rule to contradict. Parallelism is
a \code{join}, written where it happens, and a unit with one process does
not write one.

\subsection{Processes cannot take a mutable borrow}

The obvious way to write the two processes does not compile, and finding
that out cost ten lines.

\begin{lstlisting}[caption={What does not compile, and the error it gives. This is the code a Rust programmer writes first.},label={lst:e0499}]
async fn port_a(&mut self) { /* ... */ }
async fn port_b(&mut self) { /* ... */ }

// error[E0499]: cannot borrow `*self` as mutable
//               more than once at a time
join2(self.port_a(), self.port_b()).await;
\end{lstlisting}

Hardware processes share state: they drive registers. Rust permits one
mutable borrow at a time. The two facts collide directly.

The repair is that a register is a cell, and a process takes a shared
reference.

\begin{lstlisting}[caption={The repair. A register is a cell several processes reach, which is what a register is. Reading one is \code{async}; driving one is not.},label={lst:reg}]
pub struct Reg<T: Copy>(Cell<T>);

impl<T: Copy> Reg<T> {
    pub async fn get(&self) -> T { tick().await; self.0.get() }
    pub fn set(&self, x: T) { self.0.set(x) }
    pub fn peek(&self) -> T { self.0.get() }   // testbench only
}

async fn port_a(&self) {
    loop {
        let hits = self.hits_a.get().await;
        self.hits_a.set(hits + 1);
    }
}
\end{lstlisting}

\code{run} still takes \code{\&mut self}, and the processes it starts take
\code{\&self}. That compiles, and it also describes the hardware better: a
register is a thing several processes reach, which is what a register is.

\subsection{A register read is the wait, and a process loops}
\label{sec:e-regread}

Listing~\ref{lst:reg} makes two claims that earlier drafts did not, and
both follow from what a register is.

\textbf{Reading a register is \code{async}.} A wire has no clock edge, so
\code{In::get} is plain and returns at once. A register holds the value
latched at the last edge, so reading it is waiting for that edge, and
\code{Reg::get} is \code{async} to say so. Driving a register states the
next value and waits for nothing, so \code{set} is plain. The
\code{.await} on a read is the mark the lowering turns into the
register; the value alive across it is what gets stored. \code{peek}
reads without waiting, and exists for testbenches, for the reason
\code{assert!} exists: it observes and drives nothing, and synthesis
skips it.

\textbf{A process loops.} A unit runs for as long as the clock does, so
every leaf \code{run} is a \code{loop}, and each iteration contains a
register read. That read is where the iteration waits for the edge, so
one iteration is one cycle. A parent whose \code{run} only joins its
children needs no loop of its own: the children loop, the join never
completes, and the parent is continuous through them.

The prototype enforces this in the only way a simulation can. A register
read yields to the executor once, so one poll of the top unit is one
cycle, and a loop that reads no register and awaits no operator never
yields and spins. That is the right behaviour, because such a loop is a
combinational loop in the hardware, which is the thing that was wrong.

\begin{pitfall}
This is the constraint most likely to surprise somebody writing a unit for
the first time, because the code that fails is the code a Rust programmer
would write first. It should be in the library's documentation and in the
error message, not only in a paper.
\end{pitfall}


\subsection{A unit that contains units}
\label{sec:e-hierarchy}

A submodule is a field. A bus that joins two submodules is a field too,
owned by whichever unit contains both of its ends.

\begin{lstlisting}[caption={Submodules and the wire between them are fields of the parent. Each child is generic over its own configuration, not the parent's.},label={lst:hier}]
pub struct Top<TC: TopConfig> {
    pub producer: Producer<TC::P>,
    pub consumer: Consumer<TC::C>,
    pub link: DataBus,
    pub pes: [Pe; 4],
}
\end{lstlisting}

\subsection{Configurations nest}
\label{sec:e-nestedconfig}

Section~\ref{sec:e-binding} gives a design one configuration and makes
every unit generic over it. That holds while every choice is global. It
stops holding as soon as a submodule has choices of its own, and
Listing~\ref{lst:hier} is where that shows.

A configuration is a tree. A unit declares the configuration \emph{it}
needs, and a parent's configuration names its children's rather than
holding their fields.

\begin{lstlisting}[caption={Each unit declares what it needs. The parent names its children's configurations and never mentions their knobs.},label={lst:nestedcfg}]
pub trait ProducerConfig {
    type Gen: Generator;
    const DEPTH: usize;
}

pub trait ConsumerConfig {
    const DEPTH: usize;        // the same name, a different unit
    const CHECKSUM: bool;
}

pub trait TopConfig {
    type P: ProducerConfig;
    type C: ConsumerConfig;
    const NAME: &'static str;
}
\end{lstlisting}

Two things follow, and both were the reason to do it.

\textbf{Two children may want a knob of the same name.} Both
\code{DEPTH}s above exist, differ, and never meet, because each belongs
to its own trait. A flat configuration would have had to rename one of
them after the fact, which is a rename in a design that had nothing wrong
with it.

\textbf{A child's configuration is reusable.} It is written once and
named by whichever parents want it, so two builds may differ in one
subsystem and share the rest.

\begin{lstlisting}[caption={Two builds. The consumer's configuration is written once and used by both; only the producer differs.},label={lst:nestedbuilds}]
pub struct CheckedConsumer;
impl ConsumerConfig for CheckedConsumer {
    const DEPTH: usize = 64;
    const CHECKSUM: bool = true;
}

pub struct Fpga;
impl TopConfig for Fpga {
    type P = FastProducer;
    type C = CheckedConsumer;
    const NAME: &'static str = "fpga";
}

pub struct Tiny;
impl TopConfig for Tiny {
    type P = SmallProducer;
    type C = CheckedConsumer;   // reused, unchanged
    const NAME: &'static str = "tiny";
}
\end{lstlisting}

Nesting costs nothing at elaboration. A child's constant is still a
compile-time value reached through the parent, so it may still size an
array:

\begin{lstlisting}[caption={A child's constant, reached through the parent, still sizes an array. This is the long spelling.},label={lst:nestedconst}]
pub const DEPTH: usize =
    <<Fpga as TopConfig>::P as ProducerConfig>::DEPTH;
pub static BUF: [u32; DEPTH] = [0; DEPTH];
\end{lstlisting}

The spelling in Listing~\ref{lst:nestedconst} is the one real cost of
nesting: reaching two levels down needs a fully qualified path, and three
levels reads worse again. The repair is one type alias per level of the
tree, written once beside the trait, so it serves every build rather than
one.

\begin{lstlisting}[caption={The repair. A generic alias per level, and a per-build alias where one build is being talked about. Both compile in probe~17.},label={lst:nestedalias}]
pub type ProducerOf<T> = <T as TopConfig>::P;
pub type ConsumerOf<T> = <T as TopConfig>::C;

// One qualification, not two, for any build.
pub const DEPTH: usize =
    <ProducerOf<Fpga> as ProducerConfig>::DEPTH;

// None at all, where one build is meant.
pub type FpgaProducer = ProducerOf<Fpga>;
pub const D: usize = <FpgaProducer as ProducerConfig>::DEPTH;
pub static BUF: [u32; D] = [0; D];
\end{lstlisting}

A derive on the configuration trait could emit the per-level aliases,
which would make the long spelling something nobody writes.

Wiring is not a stored back-reference. A child is handed the ports it
needs when its \code{run} is called, and that is what the inputs and
outputs of \code{Module} are for.

\begin{lstlisting}[caption={Instantiating two submodules and wiring them. Both children receive the same bus; neither holds a reference to the other.},label={lst:hier-run}]
impl<TC: TopConfig> Module<(), ()> for Top<TC> {
    async fn run(&mut self, _i: (), _o: ()) {
        join2(
            self.producer.run((), ProducerOut { out: &self.link }),
            self.consumer.run(ConsumerIn { inp: &self.link }, ()),
        ).await;
    }
}
\end{lstlisting}

\subsection{Why the borrow rule does not bite here}

Section~\ref{sec:e-units} showed that two processes of one unit cannot
both take \code{\&mut self}. Listing~\ref{lst:hier-run} appears to do
exactly that and compiles, which is worth explaining rather than leaving
as an inconsistency.

The two cases are not the same. A process is a method on the unit, so it
borrows the whole of \code{self} and two of them collide. A submodule is a
\emph{field}, and Rust permits simultaneous mutable borrows of disjoint
fields, so \code{\&mut self.producer} and \code{\&mut self.consumer} coexist
legally. The bus between them is borrowed shared by both, which is what a
wire is.

\begin{rulebox}[Where interior mutability is needed]
State that two processes of one unit both drive. Not submodules, which are
disjoint fields, and not the buses between them, which are shared by
design.
\end{rulebox}

\subsection{Arrays of instances}

An array of identical submodules is an array, and joining them is a join
over an iterator.

\begin{lstlisting}[caption={An instance array. The replication count is a Rust array length, so a config can set it.},label={lst:hier-array}]
join_all(
    self.pes.iter_mut()
        .enumerate()
        .map(|(i, pe)| pe.run(i as u32, ()))
).await;
\end{lstlisting}

TxHDL wrote this as \code{instance pe[16]} with a \code{connect} loop
beside it. Here the replication is an array length and the connection is
what each element is passed, so a config may set the count through an
associated constant and nothing else changes.

% SPDX-License-Identifier: Apache-2.0
\section{Signals, Ports and Interfaces}
\label{sec:e-wiring}

These three were the muddiest part of the merged specification, and LHDL's
version of them is where its grammar and its examples disagreed most. The
embedding settles them, and the settlement is one idea.

\begin{rulebox}[Direction is not a property of a wire]
The same wire is driven at one end and read at the other. Direction
belongs to the end, not to the wire, so a unit never holds a wire. It
holds an end.
\end{rulebox}

\subsection{A signal has two ends}

Creating a signal yields its ends separately, the way
\code{std::sync::mpsc} yields a sender and a receiver.

\begin{lstlisting}[caption={A signal is created as a pair. \code{T} is a struct, so a wire may be a bundle without a second concept for the bundled case.},label={lst:split}]
pub fn signal<T: Copy + Default>() -> (Out<T>, In<T>);

let (tx, rx) = signal::<Beat>();
\end{lstlisting}

\code{Out<T>} has \code{set} and no \code{get}. \code{In<T>} has \code{get}
and no \code{set}. A unit given the reading end cannot drive, and the
compiler says so in the plainest possible terms:
\code{error[E0599]: no method named `set` found for struct `In`}.

\subsection{Cardinality is the opposite of mpsc, and the types say so}

A channel in software has many senders and one receiver. A wire has one
driver and many readers. The library inverts the derive accordingly.

\begin{table}[H]
\caption{Why \code{In} is \code{Clone} and \code{Out} is not.}
\label{tab:cardinality}
\centering
\footnotesize
\begin{tabular}{@{}lll@{}}
\toprule
 & Sender or driver & Receiver or reader \\
\midrule
\code{mpsc} & \code{Clone}     & unique \\
a wire      & unique           & \code{Clone} \\
\bottomrule
\end{tabular}
\end{table}

That one line of derive buys the oldest error message in hardware design.
A second driver on a net would be a second \code{Out}, \code{signal}
returns exactly one, and \code{Out} is not \code{Clone}, so there is no
way to make another.

\begin{lstlisting}[caption={Two drivers on one wire. The classic multiple-driver error is a move error.},label={lst:twodrivers}]
let (tx, _rx) = signal::<u32>();
let a = Driver { out: tx };
let b = Driver { out: tx };
// error[E0382]: use of moved value: `tx`
\end{lstlisting}

Nothing in the language had to define that rule. Fanout, meanwhile, is
\code{rx.clone()} and costs nothing, because reading a wire is free.

\subsection{An interface is a struct of ends}

An interface is a named group. Its members are signals and channels, and
it holds no direction, because its members do not.

A \term{role} is the pairing that says which end of each member this side
takes. A \term{port} is what a unit is handed: an interface's ends,
selected by a role.

\begin{lstlisting}[caption={An interface and the two roles over it. Each role names the end it takes for every member.},label={lst:iface}]
pub struct Wishbone {
    pub adr: Signal<U<32>>,
    pub ack: Signal<Bit>,
    pub dat: Chan<Beat>,
}

pub struct InitiatorPort {
    pub adr: Out<U<32>>,
    pub ack: In<Bit>,
    pub dat: Tx<Beat>,
}

pub struct TargetPort {
    pub adr: In<U<32>>,
    pub ack: Out<Bit>,
    pub dat: Rx<Beat>,
}
\end{lstlisting}

The two role structs are mechanical: each field is the same member with
the other end. A derive writes them, which is what the \code{modport}
keyword used to do, and it cannot get them out of step with the interface
because it reads the interface to make them.

\subsection{Where the ends go}

A unit does not store a reference to the interface. It is handed its port
when \code{run} is called, which is what the inputs and outputs of
\code{Module} are for, and what makes wiring a call rather than a graph of
back-references.

\begin{lstlisting}[caption={The parent creates the interface and hands each child the end it takes.},label={lst:wire-up}]
impl Module<(), ()> for Top {
    async fn run(&mut self, _i: (), _o: ()) {
        join2(
            self.cpu.run(self.bus.initiator(), ()),
            self.ram.run(self.bus.target(), ()),
        ).await;
    }
}
\end{lstlisting}

\begin{table*}[t]
\caption{The four concepts, and what each one is responsible for. Reading
down the table is the life of a wire.}
\label{tab:wiring}
\centering
\footnotesize
\begin{tabular}{@{}llll@{}}
\toprule
Concept & Is & Knows about direction & Held by \\
\midrule
\code{Signal<T>}, \code{Chan<T>} & one wire, or one handshaked channel & no & nobody \\
\code{Out<T>} / \code{In<T>}, \code{Tx<T>} / \code{Rx<T>} & one end of one member & yes, by which methods exist & a unit \\
interface & a named group of members & no & the unit containing both ends \\
role and port & a group of ends, one per member & yes, by construction & passed to \code{run} \\
\bottomrule
\end{tabular}
\end{table*}

\input{embedding_sections/08_binding}
% SPDX-License-Identifier: Apache-2.0
\section{The Top Is a Program}
\label{sec:e-main}

Section~\ref{sec:e-hierarchy} put units inside units. The outermost one
has to be reached from somewhere, and in an embedded language that
somewhere already exists.

\begin{lstlisting}[caption={A design is a program. The config names the top unit as well as filling it, so \code{main} needs nothing else.},label={lst:main}]
pub trait Config {
    type Top: Module<(), ()>;
    const NAME: &'static str;
    const CLK_HZ: u64;
    fn top() -> Self::Top;
}

fn main() {
    let which = std::env::args().nth(1)
        .unwrap_or_else(|| "fpga".to_string());
    match which.as_str() {
        "asic" => elaborate::<Asic>(),
        _      => elaborate::<Fpga>(),
    }
}
\end{lstlisting}

There is no separate elaboration entry point and no build description
language. Probe~13 is a binary rather than a library for that reason, and
running it prints what it elaborated:

\begin{lstlisting}[style=ascii,caption={Running the design. Each arm of the match is monomorphised separately, so both builds are compiled and either can be asked for.},label={lst:run}]
$ bazel run //experiments/rust_embedding:probe_main -- asic
elaborated config 'asic' at 400 MHz
\end{lstlisting}

\subsection{A build is one item}

Written out, a build was a unit struct, an \code{impl} of the design's
own configuration trait, and an \code{impl Config} whose \code{fn top}
constructed the entire hierarchy field by field. Most of that is
mechanical, and the constructor is the part that grows with the design.

Two changes to the runtime remove it, and a macro removes the rest.
\code{Config::top} builds the top from its \code{Default}, because a
register's default is its reset value and a socket's is its
implementation, so a unit that derives \code{Default} needs no
constructor at all. A configuration is a unit struct that derives
\code{Default} too, because a unit generic over it derives
\code{Default} and the derive asks that of every parameter. Then
\code{config!} writes the struct and both impls:

\begin{lstlisting}[caption={A build is one item. The body is spliced into the design's own configuration impl unchanged, so the macro never parses it.},label={lst:config-macro}]
config! { Fpga: TopConfig for Top<Fpga> {
    type Filter = SixteenTaps;
    const CLK_HZ: u64 = 100_000_000;
} }

config! { Asic: TopConfig for Top<Asic> {
    type Filter = SixtyFourTaps;
    const CLK_HZ: u64 = 400_000_000;
} }
\end{lstlisting}

The example in the appendix went from 81 lines to 55 under this change,
and every line removed was one that repeated a field name.

Three things follow, and the third is the one that matters.

The config gains \code{type Top} and \code{fn top()}, so it says which
design to elaborate as well as what to fill it with. A config is then the
whole of what a build is, which is what the configuration facet was for.

Selecting a build is an ordinary argument. Each arm names a different type
and is compiled separately, so asking for a build that does not exist is a
compile error rather than a missing file.

And the toolchain becomes a library rather than a program that reads
source. \code{elaborate} is a function the design calls, so a project may
wrap it, run two configs and compare them, generate a report, or drive a
simulation from a test, using whatever Rust already offers. None of that
needed designing.

% SPDX-License-Identifier: Apache-2.0
\section{What Was Compiled}
\label{sec:e-probes}

Every claim in this paper about what Rust does is a probe that was built.
The probes are one file per question in \code{experiments/rust\_embedding},
and thirteen of them are expected to fail, so they are tagged \code{manual}
and stay out of the default build. A failing probe is kept rather than deleted,
because the failure is the answer.

\begin{table*}[t]
\caption{The probes and what they answered. Rust 1.90.0 stable unless
stated. Two results changed the design rather than confirming it.}
\label{tab:probes}
\centering
\footnotesize
\begin{tabular}{@{}lp{60mm}p{80mm}@{}}
\toprule
Probe & Question & Answer \\
\midrule
\code{probe\_widths}
  & Does \code{U<\{A + B\}>} compile?
  & \textbf{No.} \code{error: generic parameters may not be used in const operations} \\
\code{probe\_widths\_nightly}
  & Does it compile on nightly?
  & \textbf{Yes}, with \code{generic\_const\_exprs}, and \code{mul(a, b)} resolves to \code{U<64>} at the call site \\
\code{probe\_design}
  & Is \code{async fn} in a trait usable?
  & \textbf{Yes}, warning that auto trait bounds such as \code{Send} cannot be stated \\
\code{probe\_design}
  & Inputs and outputs as type parameters?
  & \textbf{Yes.} One struct implements \code{Unit} twice with different shapes \\
\code{probe\_design}
  & Interface as struct, modport as trait?
  & \textbf{Yes}, including a unit generic over the role rather than the bus \\
\code{probe\_comp}
  & Do \code{Bus}, \code{Port}, \code{Signal}, \code{Chan} hold together?
  & \textbf{Yes}, and construction compiles \\
\code{probe\_comp}
  & Is a \code{Tag} type parameter workable?
  & \textbf{Yes}, with the policy as associated constants readable in the body \\
\code{probe\_attr\_block}
  & Is \code{\#[tag(..)]} on a block stable?
  & \textbf{No.} \code{error[E0658]: attributes on expressions are experimental} \\
\code{probe\_processes}
  & Can \code{run} join two processes?
  & \textbf{Yes}, with processes taking \code{\&self} and registers as cells \\
\code{probe\_processes\_mut}
  & Can they take \code{\&mut self}?
  & \textbf{No.} \code{error[E0499]: cannot borrow *self as mutable more than once at a time} \\
\code{probe\_config}
  & Is a config a struct with associated types and constants?
  & \textbf{Yes}, and its constants work in an array length, so they are genuinely compile-time \\
\code{probe\_item\_order}
  & May an \code{impl} precede the \code{struct} it is for?
  & \textbf{Yes.} Items in a module are not order dependent, including generically and for types declared later \\
\code{probe\_derive}
  & Does a derive give the same reading and stay idiomatic?
  & \textbf{Yes}, in about thirty lines needing neither \code{syn} nor \code{quote} \\
\code{probe\_funcs}
  & Can operator latency replace a hand-placed stage boundary?
  & \textbf{Yes.} An \code{async} operator in \code{crate::pipeline} is awaited, and the depth follows from the operators used \\
\code{probe\_if}
  & Do \code{when!} and \code{with!} work, the fields as entries, the condition first or the struct first?
  & \textbf{Yes}, as procedural macros; \code{macro\_rules!} cannot follow an \code{expr} with \code{?} or \code{<=} \\
\code{probe\_when\_bad}
  & An entry with no colon?
  & \textbf{Refused:} \code{expected `field: value`, `c ? field: value` or `c ? \{ .. \}`} \\
\code{probe\_case}
  & Does \code{case!} take Rust patterns, with guards and alternatives?
  & \textbf{Yes}, tested with \code{matches!}; the first arm that matches wins \\
\code{probe\_case\_bad}
  & A drive with no arrow in an arm?
  & \textbf{Refused:} \code{expected `register <= value`}, the macro's own message \\
\code{probe\_case\_binding}
  & May an arm's pattern bind, as \code{Op::Load(v)}?
  & \textbf{No.} \code{error[E0425]: cannot find value `v`}; \code{matches!} scopes its bindings \\
\code{probe\_slice}
  & May a slice bound be a const parameter, or a run-time offset?
  & \textbf{Yes} to both, with the width static \\
\code{probe\_slice\_dyn}
  & May the slice \emph{width} be a run-time value?
  & \textbf{No.} \code{error[E0435]: attempt to use a non-constant value in a constant} \\
\code{probe\_hierarchy}
  & How is a unit with submodules instantiated and wired?
  & Submodules and their buses are fields; ports are passed to \code{run}. Disjoint fields, so the borrow rule of probe 5b does not apply \\
\code{probe\_main}
  & Can \code{fn main()} instantiate the top through a config?
  & \textbf{Yes}, and it runs: \code{bazel run ... -- asic} prints \code{elaborated config 'asic' at 400 MHz} \\
\code{probe\_ports}
  & Can direction be enforced rather than documented?
  & \textbf{Yes}, by which methods an end exposes \\
\code{probe\_ports\_bad}
  & Can a unit drive an input?
  & \textbf{No.} \code{error[E0599]: no method named `set` found for struct `In`} \\
\code{probe\_values}
  & Do \code{Bit}, \code{Logic}, \code{U<N>}, \code{I<N>} work on stable?
  & \textbf{Yes}, including the IEEE 1164 resolution table and \code{U} as the default of an undriven wire \\
\code{probe\_signal\_split}
  & Can a signal yield its two ends separately, as \code{mpsc} does?
  & \textbf{Yes}, with \code{In} clonable for fanout and \code{Out} unique \\
\code{probe\_two\_drivers}
  & Can a wire have two drivers?
  & \textbf{No.} \code{error[E0382]: use of moved value} \\
\code{probe\_nested\_config}
  & Can a submodule have a configuration of its own?
  & \textbf{Yes.} A parent's config names its children's, so two children may share a knob name and a child's config is reusable \\
\code{probe\_interface\_macro}
  & Can \code{macro\_rules!} generate the role views?
  & \textbf{Yes}, for exactly two roles, spelling \code{in} as \code{inp}, and a role that omits a member compiles silently \\
\code{probe\_interface\_proc}
  & Can a procedural macro do it properly?
  & \textbf{Yes.} Any number of roles, \code{in} as written, under two hundred lines, no \code{syn} \\
\code{probe\_interface\_incomplete}
  & A role that omits a member?
  & \textbf{Refused:} \code{role `Target` does not say which end of `dat` it takes} \\
\code{probe\_interface\_twodrivers}
  & Two roles driving one member?
  & \textbf{Refused:} \code{member `adr` is driven by more than one role: A, B} \\
\code{probe\_domain}
  & Can the clock domain be a type parameter with a named default?
  & \textbf{Yes.} A single-clock design never mentions one; a two-clock design names the second \\
\code{probe\_domain\_cross}
  & A crossing without a \code{Crossing}?
  & \textbf{Refused:} \code{expected In<u32, Clk400>, found In<u32, Clk100>}, and the default is not a wildcard \\
\code{probe\_infer}
  & May a const width be left to inference, as \code{\_}?
  & \textbf{Yes}, since Rust 1.89, where the context fixes it: an annotated destination, a typed use, the other operand, the return type \\
\code{probe\_infer\_bad}
  & A width nothing fixes, a slice compared with a number?
  & \textbf{No.} \code{error[E0284]: type annotations needed}; and the lowering, which sees no types, asks for every width it needs \\
\code{probe\_infer\_lower}
  & A width Rust would infer, under \code{\#[lower]}?
  & \textbf{Refused:} \code{`slice::<..>` needs its width written: the lowering sees no types}; \code{concat}'s low operand is the one width it does not need \\
\bottomrule
\end{tabular}
\end{table*}

\subsection{The four that changed the design}

\textbf{Arbitrary-width arithmetic needs nightly.} A hardware library wants
to write that multiplying an $A$-bit value by a $B$-bit value gives an
$A+B$-bit one. Stating that in a signature needs arithmetic over const
generic parameters, and stable Rust refuses it. The choices are nightly
with \code{generic\_const\_exprs}, or one function per width pair, which
compiles on stable and does not scale. This is the largest single
constraint on the proposal, and it was a guess until it was built.

\textbf{Parallel processes cannot take a mutable borrow.}
Section~\ref{sec:e-units} gives the collision and the repair. The point
here is that ten lines of Rust found a constraint that reading the language
reference had not suggested, and that the repair, a register being a cell
several processes reach, is a better model of a register than the one that
failed.

\textbf{A \code{macro\_rules!} conditional cannot reach inside a block.}
Writing \code{with!} showed where the macro route stops, and that is the
strongest argument for the recommendation in
Section~\ref{sec:e-elaboration}. Section~\ref{sec:e-walls} states it.

\textbf{One driver per wire comes free.} It was not designed in.
\code{signal} returns one \code{Out} and \code{Out} is not \code{Clone},
so a second driver is a value that cannot be made, and the oldest error
message in hardware design arrives as \code{error[E0382]: use of moved
value}. Section~\ref{sec:e-wiring} states the cardinality this rests on.

\subsection{On method}

The first draft of this work argued from the language rules and was wrong
twice in eleven claims. Both errors were in the direction of optimism, and
both were cheap to find. A design document for an embedded language should
include its probes, and the probes should be part of the build, so that a
Rust release that changes an answer breaks something visible rather than
quietly invalidating a paragraph.

% SPDX-License-Identifier: Apache-2.0
\section{What Rust Cannot Do}
\label{sec:e-walls}

One wall is left, and it is the one a separate parser was buying.

\subsection{Control flow on a signal}

A signal is a node in a graph. It has no value while the Rust program
runs, so it cannot be a condition.

\begin{lstlisting}[caption={The wall. A signal has no value while the program runs, so it cannot be a condition.},label={lst:wall}]
if cond { a } else { b }   // cond must be a bool.
                           // A signal is not a bool.
\end{lstlisting}

\code{if} and \code{match} take a \code{bool} and a pattern, and neither
can be overloaded. Control flow on a signal needs another spelling.

\subsubsection{The name is not \code{if\_!}}

The obvious macro name is \code{if\_!}, and it should be rejected for a
better reason than the trailing underscore. \code{if} branches: one arm
runs and the other does not. This does not branch. Both arms exist in the
hardware at once and a multiplexer picks between them, so area is paid for
both and a side effect in the arm not chosen still has to be suppressed. A
reader who brings Rust's \code{if} semantics to a construct named after it
is misled about the one thing that matters most.

\code{when!} is what Chisel~\cite{bachrach2012} and SpinalHDL call it, so
the engineers who will read it already know what it means, and it claims
nothing it does not do.

\subsubsection{Two forms, and where the limit is}

An expression needs no macro at all. \code{mux} is the operation, and a
function says so.

\begin{lstlisting}[caption={The expression form is a function. There is nothing to predicate.},label={lst:mux}]
let chosen = mux(reset, Sig(0), self.count.get());
\end{lstlisting}

A statement form has to predicate whatever is written inside it, and that
is what real logic needs.

\begin{lstlisting}[caption={The statement form. It reads as a conditional and lowers to predicated drives. The arrow points into the register: the register becomes the value.},label={lst:when}]
when!(enable => {
    self.count <= count.wrapping_add(1);
    self.flag  <= Bit::One
} else {
    self.count <= 0;
    self.flag  <= Bit::Zero
});
\end{lstlisting}

Probe~10 compiles both. The arrow is \code{<=}, VHDL's signal
assignment, because it points into the register and says that the
register becomes the value. It is also the reason \code{when!} is a
procedural macro. The first version was \code{macro\_rules!} and spelled
the arrow \code{=>}, not by choice: an \code{expr} fragment may be
followed only by \code{=>}, \code{,} or \code{;}, so \code{lhs <= rhs}
cannot be matched declaratively at all. That is the third place
\code{macro\_rules!} ran out, after the two in
Section~\ref{sec:e-wiring}, and each time the reason was the same: a
declarative macro matches a grammar, and the grammar it can match is not
the one the design wants to write.

The procedural version splits each statement on the first \code{<=} at
depth zero, so the left side may be any path, and refuses a statement
that has none:

\begin{lstlisting}[style=ascii,caption={A drive without an arrow. The macro's own message, not a type error three expansions later.},label={lst:when-err}]
error: expected `register <= value`
\end{lstlisting}

What it still does not do is reach inside an arbitrary block and rewrite
the statements it finds there. That is the difference between predicating
a list of drives and reinterpreting a plain \code{if}, and it is the
elaboration decision of Section~\ref{sec:e-elaboration} rather than a
macro's.

\subsection{Bit slicing}

\code{w[0..8]} cannot work, because \code{Index} returns a reference to
something that already exists and a slice of a signal is a new node. It
becomes \code{w.slice::<0, 8>()}.

The bounds do not all have to be literals, and the rule is not the one a
first reading suggests.

\begin{table}[H]
\caption{What a slice bound may be. The width is the type, so it is the
one that cannot vary.}
\label{tab:slice}
\centering
\footnotesize
\begin{tabular}{@{}lll@{}}
\toprule
Form & Offset & Width \\
\midrule
\code{w.slice::<0, 8>()}        & literal          & literal \\
\code{w.slice::<LO, LEN>()}     & const parameter  & const parameter \\
\code{w.slice\_at::<8>(i)}      & \textbf{run time} & static \\
\code{w.slice\_var(i, n)}       & run time         & \textbf{rejected} \\
\bottomrule
\end{tabular}
\end{table}

A generic function may slice without knowing the numbers, because a const
generic parameter is a compile-time value and not a literal. An offset may
be a run-time value, which synthesises to a barrel shifter feeding a
fixed-width slice.

A width may not. The width is the type, and a type cannot depend on a
value that exists only at run time: probe~11b gives
\code{error[E0435]: attempt to use a non-constant value in a constant}.
That is not a Rust limitation to work around. A signal whose bit width is
decided at run time is not a thing hardware has, so the compiler is
refusing to express something that could not be built.

\section{What the Embedding Costs}
\label{sec:e-costs}

\subsection{Structural typing on data changes shape}

The merged specification resolved one of its four conflicts by making
transaction compatibility structural: two transactions with matching fields
are compatible whatever they are called~\cite{unification}. Rust is
nominally typed.

A derive can recover it by generating an associated layout type, so that
compatibility becomes a bound rather than a name comparison. It recovers it
positionally, though, so the rule becomes ``the same field types in the
same order'' rather than ``the same fields''. That is a real change to a
decision already taken, and unlike everything else in this paper it has not
been compiled.

\subsection{A thesis is inverted}

The project's second thesis~\cite{theses} names making the language a valid
program in an existing language as one direction, and then takes the other,
on the grounds that a full toolchain has become affordable to build. This
paper takes the direction the thesis declined.

The thesis is the rationale section of the merged specification, so it
should be rewritten rather than quietly contradicted. Its argument survives
the change of direction, and only its conclusion moves: a toolchain is
still being built, but it buys a front end over Rust's own syntax tree
rather than a lexer and a parser, and the type system, the module system,
the generics and the tooling arrive already written.

\subsection{Smaller frictions}

\code{concat!} is a \code{std} macro, so the bit concatenation macro needs
another name. \code{async fn} in a public trait warns that auto trait
bounds cannot be stated, which is harmless while elaboration is single
threaded and would matter if it ever were not.

\input{embedding_sections/12_elaboration}
% SPDX-License-Identifier: Apache-2.0
\section{Related Work}
\label{sec:e-related}

Three systems bound this design, and it is closer to each than to the
others.

\textbf{Chisel}~\cite{bachrach2012} is the structural half. A hardware
module is a class in a general-purpose language, a bundle of wires has a
direction per field through \code{Flipped}, control flow on a signal is
\code{when} and \code{Mux}, and elaboration is the program running and
recording a graph. Sections~\ref{sec:e-units} to~\ref{sec:e-binding} are
that architecture in Rust, and they should be read as a port rather than
a proposal. SpinalHDL follows the same shape in the same host language.
Two things differ, both because of the host. Rust's ownership makes a
second driver on a wire a move error at compile time, where Chisel
detects it at elaboration. And Rust's trait system lets a configuration be
a type, so a build is a type alias and a binding path cannot go stale.

\textbf{Calyx}~\cite{calyx} is the nearest ancestor of the behavioural
half, and it is an intermediate representation rather than a language.
It separates structure, the wires and groups that compute, from control,
a program that says when groups run. Its control has \code{seq}, which
runs groups one after another, and \code{par}, which runs them at once.
Section~\ref{sec:e-units} arrives at the same two operators from the other
direction: an \code{await} is a \code{seq} boundary, and \code{join!} is
\code{par}. Calyx lowers both to a state machine driving the groups, and
that lowering is the one this paper has not yet done.

\textbf{Filament}~\cite{filament} types each signal with the interval of
cycles over which it is valid, so that composing two components with
different latencies is checked rather than simulated. That is what an
\code{async} operator whose latency the mapping decides is reaching for.
Filament makes the latency part of the type and checks it statically;
this paper defers it to the mapping and would learn it at elaboration.
Filament's is the stronger guarantee, and the price is that latency has
to be stated in the type. Whether an \code{async fn} can be given a
timeline type after the fact is an open question.

\textbf{RHDL}~\cite{rhdl} is the nearest thing in Rust, and the source
of one decision here. It embeds a hardware description language in Rust
with the clock domain as a type parameter of the signal, so that a
crossing without a synchroniser is a type error. Section~\ref{sec:e-wiring}
takes that directly, and departs from it only in what the default domain
means. RHDL's structural model is Chisel's; its behavioural model is
Rust functions on values, evaluated to produce combinational logic, with
registers explicit. It has no \code{async}, so the sequencing half of
this paper is as far from RHDL as from Chisel.

\textbf{XLS}~\cite{xls} is the one to weigh most carefully, because it
is a working toolchain that has done what Section~\ref{sec:e-next} calls
the next experiment. Its front end, DSLX, is a dataflow language with
Rust's surface syntax: \code{fn}, \code{let}, \code{match}, structs,
enums and parametric types, read by a parser of its own. A DSLX function
is scheduled into a pipeline by the compiler, with the stage boundaries
chosen to meet a clock period rather than placed by the designer, which
is what Section~\ref{sec:e-pipeline} asks of an \code{async fn}. A DSLX
\code{proc} is a process with state and channels, sending and receiving
in a loop, which is the sequenced half of this paper under another
name. XLS emits Verilog, and it has been used in silicon.

That makes XLS the nearest ancestor of the behavioural half, nearer than
Calyx, and it also makes it the counterexample to this paper. XLS is the
path the project's second thesis~\cite{theses} chose: a new language,
shaped like Rust, with its own toolchain. The merged specification of
the companion paper~\cite{merge} is on that path. This paper leaves it,
and XLS is the standing evidence that the path it leaves works. The
argument for leaving it is therefore not that a separate language cannot
be made to lower; XLS shows it can. It is that the type system, the
module system, the generics, the package manager and the editor arrive
already written when the host is Rust itself, and that a DSLX-shaped
language has to rebuild each of them. Whether that trade is worth the
one wall of Section~\ref{sec:e-walls} is the question the next
experiment answers.

What is not near. Bluespec~\cite{nikhil2004} schedules guarded atomic
rules, which is a different model of concurrency from processes that
await. High-level synthesis from C or C++ starts from sequential code and
infers a pipeline, which is the same ambition as
Section~\ref{sec:e-pipeline}, from a language that has no way to say
where a boundary is allowed to fall.

Prior work on lowering coroutines or \code{async} functions to hardware
exists and has not been surveyed for this paper. That is a gap, and it is
the first thing to close before Section~\ref{sec:e-next} is attempted,
because somebody may already have found where it breaks.

\section{What Comes Next}
\label{sec:e-next}

Nineteen probes prove that the design type-checks. Chisel is a working
compiler with a decade of use. The gap between those two facts is the
whole of what remains, and it should be closed in one place rather than
widened with more probes.

\subsection{The experiment}

Lower one function to a netlist.

\begin{lstlisting}[caption={The function to lower. Two operators, two awaits, one value live across a boundary.},label={lst:next-mac}]
async fn mac(a: U<32>, b: U<32>, prev: U<64>) -> U<64> {
    let p = mul(a, b).await;
    add(prev, p).await
}
\end{lstlisting}

The lowering has to produce, from that source and nothing else, a
two-stage pipeline in Verilog: a multiplier, a register holding \code{p}
across the boundary, an adder, and a valid signal that follows the data
through. It has to do so twice, once with \code{mul} bound to a
three-cycle implementation and once to a single-cycle one, and the second
netlist has to have one register fewer than the first without any edit to
Listing~\ref{lst:next-mac}.

That is the whole claim of Section~\ref{sec:e-pipeline} stated as a
test. If it passes, the behavioural half is real. If it cannot be made to
pass, the paper has rebuilt Chisel in Rust with better error messages,
which is worth having and is not what the project set out to do.

\subsection{What it needs}

A procedural macro on the function that reads its body as a syntax tree,
finds the awaits, and splits the body at each one into a stage. The
locals live across each split are the registers. The operators awaited
are instances whose latency is looked up in the configuration. The result
is a graph, and emitting Verilog from a graph is the part Chisel has
already shown to be routine.

Three things are known to be hard, and the experiment should be designed
to meet them rather than to avoid them.

\begin{itemize}
  \item \term{Control flow across an await.} A \code{when!} whose arms
        contain awaits is a pipeline whose stages are conditional. Calyx
        handles this with \code{if} in the control program; the lowering
        needs an equivalent, and it is the first place the async model
        might not fit.
  \item \term{Resource sharing.} Section~\ref{sec:e-pipeline} claims that
        concurrent invocations at the same await index share one
        operator. The lowering has to actually do that, or the claim is
        withdrawn.
  \item \term{Backpressure.} A stage that cannot proceed is an await that
        has not completed. The valid signal has to run backwards as well
        as forwards for that to be true in hardware, which is the elastic
        pipeline of Carloni~\cite{carloni2001}, and it has to be generated
        rather than assumed.
\end{itemize}

Since this section was first written the runtime has gained a
waveform: a probe registry named from the fields of units and a VCD
writer, so a design is watched in Surfer rather than in printed
traces. FST is the same probes with another writer, and the first
external crate the project would take. The same walk now writes a
structural netlist as well, a Verilog skeleton with modules,
registers, ports and wires and no behaviour; it sees fields only, so a
unit that takes its ports as parameters of \code{run} shows its
registers and nothing else. That gap is the argument for the
proc-macro route restated from the other side: the structure is
already recoverable at run time, and the behaviour is not.

\subsection{What it does not need}

A parser, a type checker, a module system, generics, a package manager or
an editor. That list is the argument for embedding, and it is why the
experiment is a macro of a few hundred lines rather than a compiler.

% SPDX-License-Identifier: Apache-2.0
\section{Conclusion}
\label{sec:e-conclusion}

A companion paper merged two hardware description languages and gave the
result a syntax shaped after Rust. This paper removed every construct Rust
already provides and asked what was left.

Seventeen of twenty constructs go. What remains is five Rust modules
of traits and types and a crate of macros, not a grammar.

Two deletions are worth more than the count suggests. The rule that a
timeless body may not count cycles needs no enforcement, because a cycle
count is spelled \code{.await} and a plain function has none: a paragraph
of specification and a compiler check become a fact about the type system.
And a pipeline, a sequence and a set of parallel processes turn out to be
one construct driven three ways, so \code{pipeline}, \code{stage} and
\code{pipe} all go, and the live-range analysis the last of them existed to
feed is what an async state machine does already.

Two results were negative and were worth more still, because neither was
visible from the language reference. Arbitrary-width arithmetic does not
work on stable Rust. A unit's parallel processes cannot take a mutable
borrow of the unit, and the repair, a register being a cell that several
processes reach, describes a register better than the version that failed.

One wall stands. A signal has no value while the program runs, so it cannot
be the condition of an \code{if}. That is what a separate grammar was
buying, and a macro front end over Rust's own syntax tree was the way
to buy it back; \code{when!}, \code{case!} and \code{select!} are
that front end, and \code{\#[lower]} reads them.

Half of what this paper builds is Chisel, in Rust, and the paper says so.
The structural half was settled in 2012, and rebuilding it buys better
error messages and a configuration that is a type. Those are worth having
and would not justify a new toolchain on their own.

The other half is the bet. An \code{async fn} that is a sequence, a
pipeline or a process according to how it is driven, with operator
latency the mapping's decision, is not something Chisel can express, and
it is what both source languages existed to provide. Nineteen probes
show that Rust accepts it, and the experiment of
Section~\ref{sec:e-next} shows that it lowers: an \code{async fn} to
a pipeline whose depth follows its operators' latencies, a unit to a
module, and a RISC-V core~\cite{vreteno} to a netlist that simulates
against the trace of its own source, and runs at 100~MHz on an
Artix-7. What the lowering does not read yet is stated there, and it
is the next thing to do.

The project's second thesis declined embedding on the grounds that a full
toolchain had become affordable. That argument has not weakened. It now
points the other way: if a toolchain is affordable, build the one that
arrives with a type system, a module system, generics and a package manager
already written, and spend the effort on the part nobody has built.


\begin{thebibliography}{10}

\bibitem{merge}
F.~Filmar and D.~Jankovic, ``TxHDL: Unifying a dataflow and a transaction
hardware description language,'' project repository \code{HDL/txhdl},
\code{docs/article.tex}, 2026.

\bibitem{unification}
``Unifying LHDL and TxHDL,'' project repository \code{HDL/txhdl},
\code{docs/unification-analysis.md}, 2026. States what each source language
contributes and resolves the four conflicts of meaning.

\bibitem{syntax}
``Surface syntax for the merged language,'' project repository
\code{HDL/txhdl}, \code{docs/syntax-decisions.md}, 2026. Settles the
spelling of every construct, biased toward Rust.

\bibitem{theses}
F.~Filmar, ``Some theses for a modern hardware definition language,''
project repository \code{HDL/txhdl}, \code{filmil/theses.md}, 2026.

\bibitem{embedding}
``Embedding TxHDL in Rust,'' project repository \code{HDL/txhdl},
\code{docs/rust-embedding.md}, 2026. The working notes this article is
drawn from, with the probe results in full.

\bibitem{build}
``Building the specification documents,'' project repository
\code{HDL/txhdl}, \code{docs/document-build-plan.md}, 2026.

\bibitem{bachrach2012}
J.~Bachrach \emph{et al.}, ``Chisel: Constructing hardware in a Scala
embedded language,'' in \emph{Proc. Design Automation Conf.}, 2012, pp.
1216--1225.

\bibitem{calyx}
R.~Nigam, S.~Thomas, Z.~Li, and A.~Sampson, ``A compiler infrastructure
for accelerator generators,'' in \emph{Proc. ASPLOS}, 2021, pp. 804--817.

\bibitem{filament}
R.~Nigam, P.~Kumar, and A.~Sampson, ``Modular hardware design with
timeline types,'' \emph{Proc. ACM Program. Lang.}, vol.~7, no. PLDI, 2023.

\bibitem{rhdl}
S.~Basu, ``RHDL: a hardware description language embedded in Rust,''
software, \code{https://github.com/samitbasu/rhdl}, 2023 onwards.

\bibitem{nikhil2004}
R.~Nikhil, ``Bluespec System Verilog: efficient, correct RTL from high level
specifications,'' in \emph{Proc. MEMOCODE}, 2004, pp. 69--70.

\bibitem{carloni2001}
L.~P. Carloni, K.~L. McMillan, and A.~L. Sangiovanni-Vincentelli, ``Theory
of latency-insensitive design,'' \emph{IEEE Trans. Computer-Aided Design},
vol.~20, no.~9, pp. 1059--1076, 2001.

\bibitem{rust}
N.~D. Matsakis and F.~S. Klock, ``The Rust language,'' \emph{Ada Letters},
vol.~34, no.~3, pp. 103--104, 2014.

\end{thebibliography}

\end{document}
